\section{Журнал выполнения}
\begin{enumerate}
    \item Код чтения индекса был внутри утилиты \texttt{lookup} ЛР7. Он вынесен в общий заголовок \texttt{lab7/src/index\_reader.h}, которым пользуются и \texttt{lookup}, и поиск: формат читается в одном месте. После переноса автотесты ЛР7 снова пройдены.
    \item Веб-сервер на C++ потребовал бы много кода, не связанного с поиском (сокеты, разбор HTTP). Поэтому веб-часть --- обвязка на Python, а поиск остаётся в C++. Запускать \texttt{search} на каждый HTTP-запрос значило бы каждый раз читать 1,6~ГиБ индекса (0,9~с), поэтому процесс один и постоянный, а смещение и лимит передаются в каждой строке (режим \texttt{-p}).
    \item Запрос \texttt{a \&\& b || c} по приоритетам --- это \texttt{(a \&\& b) || c}, но пользователь мог иметь в виду и \texttt{a \&\& (b || c)}. Поэтому разобранный запрос печатается со скобками вокруг каждой вложенной операции, и веб-сервис показывает его над выдачей.
    \item Слова с дефисом и апострофом: токенизатор ЛР3 делит их на части, и в индексе нет терма \texttt{hewlett-packard}. Слово запроса делится тем же токенизатором, и части соединяются через И, иначе такие запросы ничего бы не находили.
    \item Для скриншота веб-сервиса Firefox в режиме без окна отказывался работать, пока был открыт обычный Firefox с тем же профилем. Скрипт запускает его с отдельным временным профилем (\texttt{--no-remote --profile}).
    \item Корпус заменён на 1\,000\,000 документов (ЛР1), код поиска не менялся. Время запросов выросло в десятки раз вместе с длиной списков. Для замеров выдача теперь печатается по одному результату (\texttt{-n 1}): время запроса считается без вывода и от этого не зависит, а полная выдача 24 запросов по миллиону --- 6 млн строк.
    \item Первый замер на миллионе дал время ровно в 2 раза больше на всех 24 запросах: процесс, видимо, попал на энергоэффективные ядра гибридного процессора. Повторный замер совпал (в пределах 2\,\%) с замером тем же поиском в ЛР11; в отчёте --- он.
\end{enumerate}

\section{План тестирования}
Корректность выдачи проверялась на трёх уровнях.
\begin{enumerate}
    \item \textbf{Автотесты.} Тесты --- юнит-тесты на C++ (\texttt{tests/test\_search.cpp}) на собственном минимальном каркасе \texttt{lab3/tests/unittest.h} (макросы \texttt{TEST} и \texttt{CHECK}, без сторонних библиотек). Они вызывают функции программы напрямую, а последний тест запускает утилиты командной строки как отдельные процессы. Сборка и запуск --- \texttt{make test}. Эталонные значения в тестах считаются независимо от проверяемого кода, для этого используются контейнеры STL (\texttt{std::map}, \texttt{std::set}) --- это код тестов, а не поисковой системы. Веб-сервис написан на Python, поэтому его проверяет отдельный тест на Python (\texttt{tests/test\_web.py}): главная страница с формой; выдача --- 50 результатов, форма с запросом, ссылка на следующие 50; последняя страница без ссылки дальше; экранирование \texttt{<script>} из запроса; нечисловой номер страницы; 404.
\end{enumerate}

\begin{center}
\small
\begin{tabular}{|c|p{0.3\textwidth}|p{0.58\textwidth}|}
\hline
\textnumero & Тест & Что проверяется \\
\hline
1 & Операции над списками & Пересечение, объединение и разность: общий случай, пустой список, непересекающиеся списки, список с самим собой. \\
\hline
2 & Разбор запроса & Приоритеты (НЕ > И > ИЛИ), скобки, слияние вложенных однотипных операций, слово с дефисом; канонический вид разобранного запроса. \\
\hline
3 & Операторы & 5 документов (\texttt{red car}, \texttt{yellow car}, \ldots), 12 запросов с известными ответами: И, ИЛИ, НЕ, скобки, регистр, отсутствующий терм, \texttt{!missing} --- все документы. \\
\hline
4 & Толерантность разбора & 15 запросов с лишними и пропущенными пробелами, одиночными операторами, незакрытыми и лишними скобками, висящими операторами, \texttt{!!}, пустыми запросами, словом с дефисом: проверяются и выдача, и разобранный вид. \\
\hline
5 & Страницы выдачи & Границы страницы (1--50, 101--120, лимит 0 --- все, смещение за концом --- пусто); разбор строки веб-сервиса \enquote{смещение, лимит, запрос}, в т.\,ч. испорченной. \\
\hline
6 & Случайные запросы & 300 документов из 12 слов, 300 случайных запросов глубиной до 3 со всеми операторами, разными написаниями (\texttt{\&\&}, \texttt{\&}, пробел, \texttt{||}, \texttt{|}), скобками и регистром. Ответ сравнивается с вычислением того же выражения над множествами (\texttt{std::set}). \\
\hline
7 & Утилита \texttt{build/search} & Формат блока выдачи, 50 результатов, \texttt{-o}, режим \texttt{-p}, заголовок и URL в UTF-8; коды возврата: нет индекса --- 1, нет аргументов или нечисловой лимит --- 2. \\
\hline
\end{tabular}
\end{center}
\begin{enumerate}
    \setcounter{enumi}{1}
    \item \textbf{Проверки на индексе корпуса} (\texttt{scripts/bench.py}, \texttt{results/consistency.txt}) --- для 5 пар термов разной частоты проверяются тождества булевой алгебры: \texttt{N(a b)} = \texttt{N(a \&\& b)}; \texttt{N(a \&\& b)} $\le$ \texttt{min(N(a), N(b))}; \texttt{N(a || b)} = \texttt{N(a) + N(b) - N(a \&\& b)}; \texttt{N(a) + N(!a)} = 1\,000\,000; закон де Моргана \texttt{N(!(a || b))} = \texttt{N(!a !b)}; \texttt{N(a !b)} = \texttt{N(a) - N(a \&\& b)}.
    \item \textbf{Сверка с ЛР7}: число документов по одному терму совпадает с df, который выдаёт утилита \texttt{lookup} ЛР7.
\end{enumerate}

\section{Результаты}
\subsection{Автотесты и проверки}
Пройдены все 7 юнит-тестов и тест веб-сервиса. Журналы прогона (\texttt{results/tests.log} и \texttt{results/tests-web.log}):

{\scriptsize
\VerbatimInput[breaklines, breakanywhere]{../results/tests.log}

\VerbatimInput[breaklines, breakanywhere]{../results/tests-web.log}
}

Все 35 проверок на индексе корпуса пройдены (\texttt{results/consistency.txt}):

{\scriptsize
\VerbatimInput[breaklines, breakanywhere]{../results/consistency.txt}
}

\subsection{Примеры выдачи}
Дамп выдачи утилиты по запросам из \texttt{data/queries.txt} --- \texttt{results/examples.txt} (по 10 результатов). Например, запрос из задания \enquote{(красный || желтый) автомобиль} на нашем корпусе:

{\scriptsize
\begin{verbatim}
query: (red || yellow) car
parsed: (red || yellow) && car
found: 20971
1    0    's-Hertogenbosch    https://en.wikivoyage.org/wiki/%27s-Hertogenbosch
2    4    Aalborg             https://en.wikivoyage.org/wiki/Aalborg
3    5    Aalst               https://en.wikivoyage.org/wiki/Aalst
\end{verbatim}
}
а запрос с лишними пробелами по краям и между словами и без пробела после \texttt{\&\&} (\texttt{Paris \&\&museum || louvre}) разобран как \texttt{(paris \&\& museum) || louvre}, найдено 15\,375 документов.

\subsection{Скорость выполнения запросов}
Время --- разбор и выполнение запроса без вывода, внутри программы: каждый запрос выполняется 50 раз подряд, берётся среднее, из 3 запусков --- минимум (\texttt{results/bench.tsv}). Загрузка индекса (1,6~ГиБ) --- 0,9~с один раз при запуске.

{\small
\begin{longtable}{|r|p{0.6\textwidth}|r|r|}
\hline
\textnumero & Запрос после разбора & Найдено & мкс \\
\hline
\endhead
1 & \texttt{lviv} & 1\,508 & 0,6 \\
\hline
2 & \texttt{museum} & 124\,133 & 29,0 \\
\hline
3 & \texttt{the} & 987\,676 & 224,0 \\
\hline
4 & \texttt{moscow \&\& aviation \&\& institute} & 475 & 131,9 \\
\hline
5 & \texttt{cheap \&\& hostel \&\& paris} & 496 & 230,8 \\
\hline
6 & \texttt{beach \&\& surfing \&\& australia} & 672 & 118,9 \\
\hline
7 & \texttt{the \&\& and \&\& of \&\& in \&\& a} & 905\,478 & 10\,515,5 \\
\hline
8 & \texttt{ski \textbar{}\textbar{} snowboard} & 13\,941 & 13,6 \\
\hline
9 & \texttt{castle \textbar{}\textbar{} palace \textbar{}\textbar{} fortress} & 90\,555 & 501,2 \\
\hline
10 & \texttt{the \textbar{}\textbar{} and \textbar{}\textbar{} of \textbar{}\textbar{} in \textbar{}\textbar{} a \textbar{}\textbar{} to} & 999\,794 & 7\,033,0 \\
\hline
11 & \texttt{museum \&\& !art} & 75\,980 & 799,0 \\
\hline
12 & \texttt{beach \&\& !sand \&\& !swimming} & 62\,268 & 503,7 \\
\hline
13 & \texttt{!lviv} & 998\,492 & 722,2 \\
\hline
14 & \texttt{!the} & 12\,324 & 982,9 \\
\hline
15 & \texttt{!(the \textbar{}\textbar{} and)} & 2\,539 & 5\,424,3 \\
\hline
16 & \texttt{(red \textbar{}\textbar{} yellow) \&\& car} & 20\,971 & 913,4 \\
\hline
17 & \texttt{hands \&\& !feet} & 34\,569 & 280,9 \\
\hline
18 & \texttt{(hotel \textbar{}\textbar{} hostel \textbar{}\textbar{} motel) \&\& (cheap \textbar{}\textbar{} budget) \&\& !luxury} & 21\,960 & 1\,465,2 \\
\hline
19 & \texttt{(train \textbar{}\textbar{} bus \textbar{}\textbar{} ferry) \&\& (ticket \textbar{}\textbar{} pass) \&\& !airport} & 24\,660 & 2\,548,4 \\
\hline
20 & \texttt{((castle \textbar{}\textbar{} palace) \&\& (germany \textbar{}\textbar{} austria)) \textbar{}\textbar{} (temple \&\& japan)} & 21\,299 & 1\,519,0 \\
\hline
21 & \texttt{(the \textbar{}\textbar{} a) \&\& (of \textbar{}\textbar{} in) \&\& !zzzz} & 992\,588 & 4\,923,7 \\
\hline
22 & \texttt{city \textbar{}\textbar{} town \textbar{}\textbar{} village \textbar{}\textbar{} hotel \textbar{}\textbar{} museum \textbar{}\textbar{} church \textbar{}\textbar{} park \textbar{}\textbar{} river \textbar{}\textbar{} beach \textbar{}\textbar{} mountain \textbar{}\textbar{} lake \textbar{}\textbar{} island \textbar{}\textbar{} road \textbar{}\textbar{} train \textbar{}\textbar{} bus \textbar{}\textbar{} airport \textbar{}\textbar{} restaurant \textbar{}\textbar{} food \textbar{}\textbar{} market \textbar{}\textbar{} street} & 790\,259 & 23\,513,7 \\
\hline
23 & \texttt{(paris \&\& museum) \textbar{}\textbar{} louvre} & 15\,375 & 474,8 \\
\hline
24 & \texttt{lviv \textbar{}\textbar{} kyiv} & 3\,824 & 4,1 \\
\hline
\end{longtable}

}

Среднее время по 24 запросам --- 2,6~мс, медиана --- 0,61~мс, максимум --- 23,5~мс. То есть одно ядро выполняет в среднем около 400 запросов в секунду; на корпусе в 30\,000 документов было 91~мкс в среднем --- время растёт вместе с длиной списков, примерно пропорционально размеру корпуса.

\subsection{Долгие запросы}
Время запроса пропорционально суммарной длине списков, которые приходится скопировать и слить, а не числу найденных документов.
\begin{itemize}
    \item Длинное ИЛИ частых слов (запрос 22, 20 слов, 23,5~мс) --- самый долгий: объединение идёт по очереди, и накопленный результат (до 790 тыс. документов) заново сливается с каждым следующим списком, т.\,е. работа растёт как \enquote{число слов $\times$ размер результата}.
    \item И и ИЛИ самых частых слов (запросы 7, 10, 21: \texttt{the}, \texttt{and}, \texttt{of} \ldots, 4,9--10,5~мс): каждый список почти на весь миллион документов копируется и сливается целиком.
    \item НЕ (запросы 13--15, 0,7--5,4~мс): дополнение строит список из всего миллиона документов, даже для редкого терма (\texttt{!lviv} --- 998\,492 документа, 0,7~мс).
    \item Редкие термы --- доли микросекунды (\texttt{lviv} --- 0,6~мкс). И с редким термом (\texttt{moscow aviation institute} --- 132~мкс) дольше: списки всех операндов сначала целиком копируются из индекса, а пересечение начинается уже с самого короткого.
\end{itemize}
Как ускорить: объединять много списков сразу (слияние $k$ списков через кучу) или через битовую карту на $N$ битов (для миллиона документов --- 122~КиБ), начиная с коротких списков; не копировать списки документов из индекса, а проходить по ним на месте; для пересечения длинного списка с коротким --- прыжки по индексу (отдельная ЛР) или двоичный поиск; дополнение не строить явно, а хранить узел как \enquote{все, кроме списка}.

\section{Выводы}
Сделан булев поиск над индексом ЛР7: запросы с операторами И (пробел, \texttt{\&\&}), ИЛИ (\texttt{||}), НЕ (\texttt{!}) и скобками, разбор рекурсивным спуском, устойчивый к пробелам, одиночным операторам, незакрытым скобкам и висящим операторам. Есть утилита командной строки (запрос на строку, блоки выдачи) и веб-сервис из двух страниц с выдачей по 50 результатов. Корректность проверена автотестами (в т.\,ч. 300 случайных запросов против вычисления на множествах) и тождествами булевой алгебры на индексе корпуса. На корпусе в 1\,000\,000 документов запрос выполняется в среднем за 2,6~мс (медиана --- 0,61~мс), самый долгий из проверенных --- ИЛИ двадцати частых слов --- за 23,5~мс.

Недостатки: выдача упорядочена по номеру документа, а не по релевантности, поэтому первыми всегда идут статьи Wikivoyage (docID 0--34\,817), а посты Reddit --- последними: по запросу \texttt{lviv} (1\,508 документов) статья Wikipedia \enquote{Lviv} стоит на 101-м месте, \enquote{Lviv Theatre of Opera and Ballet} --- на 441-м, после статей, где Львов только упоминается, а вопросы Travel SE и посты Reddit о Львове --- ещё дальше (нужно ранжирование, например TF-IDF); нет поиска фраз (нужен координатный индекс); длинные ИЛИ и НЕ выполняются неоптимально (см. выше); индекс целиком загружается в память; веб-сервис обрабатывает запросы по одному через один процесс поиска.
